% =====================================================================
%  현장연구보고서 (Field Research Report)
%  스마트폰 기반 얼굴·음성 디지털 바이오마커를 이용한 PCOS 예측 연구
%  Author: 신홍철 (Hongchul Shin), KIST
%
%  컴파일 방법 (Korean): XeLaTeX 또는 LuaLaTeX 권장
%     xelatex main.tex   (2회 실행 — 목차/참조 갱신)
%  pdfLaTeX 사용 시에도 kotex가 대부분 처리하나, XeLaTeX를 권장합니다.
% =====================================================================
\documentclass[11pt,a4paper]{article}

% ---------- 한글 ----------
\usepackage{kotex}            % 한국어 조판 (XeLaTeX/LuaLaTeX/pdfLaTeX 호환)

% ---------- 레이아웃 ----------
\usepackage[margin=2.5cm]{geometry}
\usepackage{setspace}
\onehalfspacing

% ---------- 표 ----------
\usepackage{booktabs}
\usepackage{tabularx}
\usepackage{longtable}
\usepackage{array}
\usepackage{ragged2e}
\newcolumntype{L}{>{\RaggedRight\arraybackslash}X}

% ---------- 수식 기호 ----------
\usepackage{amsmath,amssymb}

% ---------- 링크 ----------
\usepackage[hidelinks]{hyperref}
\hypersetup{
  colorlinks=true,
  linkcolor=black,
  citecolor=blue!50!black,
  urlcolor=blue!50!black,
  pdftitle={스마트폰 기반 얼굴 음성 디지털 바이오마커를 이용한 PCOS 예측 연구 현장연구보고서},
  pdfauthor={Hongchul Shin}
}

% ---------- 제목 형식 ----------
\usepackage{titlesec}
\usepackage{abstract}

% ---------- 메타데이터 ----------
\title{\vspace{-1.0cm}\bfseries
스마트폰 기반 얼굴·음성 디지털 바이오마커를 이용한\\
다낭성난소증후군(PCOS) 예측 연구의 준비 단계 현장연구보고서\\[0.3em]
\large 선행연구 조사 · 후보 바이오마커 도출 · 데이터 수집 프로토콜 설계 및\\
\large 다기관 현장 협의 경과\\[0.5em]
\normalsize\itshape A Field Research Report on the Preparatory Phase of a Study Predicting\\
\normalsize\itshape Polycystic Ovary Syndrome (PCOS) via Smartphone-Based Facial and Voice Digital Biomarkers}

\author{신홍철 (Hongchul Shin)\\
\small 한국과학기술연구원 (Korea Institute of Science and Technology, KIST)}

\date{2026년 6월 16일}

% =====================================================================
\begin{document}

\maketitle

% ---------------- 초록 ----------------
\begin{abstract}
\noindent
\textbf{배경:} 다낭성난소증후군(polycystic ovary syndrome, PCOS)과 자궁내막증(endometriosis)은 난임의 주요 원인이나 조기 진단이 어렵고 진단 접근성·비용 문제가 크다. 본 과제는 사용자가 일상에서 스마트폰으로 자가수집한 얼굴 영상·음성·설문 데이터를 AI로 분석하여 PCOS를 \emph{스크리닝(병원 진료 권고 수준)} 하는 디지털 바이오마커를 발굴하는 것을 목표로 한다.

\smallskip
\noindent
\textbf{목적:} 본 보고서는 본격 임상연구 착수에 앞서 수행한 \emph{준비 단계의 활동 경과}를 기록한다. (1)~얼굴·음성 디지털 바이오마커 기반 질병 예측 선행연구의 체계적 조사, (2)~연구 목적에 부합하는 자가수집 후보 바이오마커 도출, (3)~데이터 수집 프로토콜(통제 환경/일상 환경) 설계, (4)~앱 개발사·임상 자문 교수와의 다기관 현장 협의 경과를 다룬다.

\smallskip
\noindent
\textbf{방법:} 2026년 4{--}6월에 걸쳐 PubMed·OpenAlex·arXiv·JMIR·Frontiers 및 국내 학술 DB(RISS·KCI·DBpia 등)를 대상으로 디지털 바이오마커 문헌을 다단계 병렬 탐색하였다. 탐색 결과는 바이오마커 카탈로그·연구설계서·수집 프로토콜로 합성하였으며, 인용 문헌은 웹 검색 기반 실존성 검증(할루시네이션 검사)을 거쳤다. 병행하여 KIST·해피문데이·고려대병원이 참여한 현장 협의 3회를 수행하였다.

\smallskip
\noindent
\textbf{결과:} 스마트폰 자가수집이 가능한 PCOS 표현형 바이오마커 9종을 A/B/C 등급으로 분류하였다(A등급: 흑색가시세포증·여드름·다모증). 또한 통제 불가 일상 환경(In-the-Wild)에서 노이즈를 제거하지 않고 측정·기록·보정하는 수집 패러다임을 설계하였다. 현장 협의 결과 수집 모달리티(얼굴·음성·설문)와 목표 표본 규모(약 900{--}1{,}500명, 초기 100명 단계적 확장)에 대한 합의가 이루어졌으나, 예산·얼굴 데이터 거부감·두 질환의 동일 파라미터 적용 타당성·IRB 재심의·기관 간 역할·지식재산권 등은 미결 쟁점으로 남았다.

\smallskip
\noindent
\textbf{결론:} 본 단계는 확정적 결론을 도출하기 위한 것이 아니라, \emph{어떤 데이터를 어떤 방법으로 수집할지 결정하기 위한 근거 기반 의사결정 과정}이었다. PCOS 중심·얼굴/음성/설문 멀티모달·일상 환경 자가수집·단계적 표본 확장을 골자로 하는 연구 방향이 수렴되고 있다.

\smallskip
\noindent
\textbf{중심어:} 디지털 바이오마커, 다낭성난소증후군, 자궁내막증, 스마트폰 자가수집, 원격 임상연구, 얼굴 영상 분석, 음성 분석, 현장연구
\end{abstract}

% ---------------- 보고서 개요 표 ----------------
\begin{table}[h]
\centering
\footnotesize
\renewcommand{\arraystretch}{1.25}
\begin{tabularx}{\textwidth}{@{}>{\bfseries}p{2.4cm} L@{}}
\toprule
보고서 유형 & 현장연구보고서 (Field Research Report) --- 연구 준비 단계 \\
작성일 & 2026-06-16 \\
연구 과제 & AI 기반 난임 예측·진단용 다중 바이오마커 패널 개발 (PCOS·자궁내막증 디지털 바이오마커 발굴) \\
수행 기관 & 한국과학기술연구원(KIST) \\
협력 기관 & 고려대학교병원(산부인과), 삼성서울병원, ㈜해피문데이(앱 개발사) \\
보고 범위 & 2026-04-27 $\sim$ 2026-06-01 현장 협의 3회 및 동 기간 수행한 문헌 조사·프로토콜 설계 \\
\bottomrule
\end{tabularx}
\end{table}

\clearpage

% =====================================================================
\section{서론}

\subsection{연구 배경}
세계보건기구(WHO) 기준 난임은 12개월 이상의 규칙적 성관계에도 임신에 이르지 못하는 상태로 정의되며, 그 주요 기질적 원인 중 하나가 PCOS와 자궁내막증이다. 두 질환은 다음과 같은 공통적 진단 한계를 가진다.

\begin{itemize}
  \item \textbf{조기 진단의 어려움:} 증상이 비특이적이고 발현이 점진적이어서 환자가 질환을 인지하기까지 오랜 시간이 걸린다.
  \item \textbf{접근성·비용 문제:} 확진을 위해서는 호르몬 검사, 초음파, (자궁내막증의 경우) 복강경 등 침습적·고비용 검사가 필요하다.
  \item \textbf{자궁내막증의 근거 부족:} 디지털 바이오마커와의 연관성에 대한 의학적 근거가 PCOS 대비 현저히 부족하다.
\end{itemize}

이러한 한계를 극복하기 위해, 본 과제는 \textbf{일상에서 비침습적으로 수집 가능한 디지털 바이오마커}(얼굴 영상·음성·앱 설문)를 AI로 분석하여, 고비용 확진 검사 이전 단계에서 질환 위험군을 스크리닝하는 접근을 제안한다. 이는 ``AI 기반 난임 예측·진단용 다중 바이오마커 패널 개발''이라는 과제 기술개발목표의 첫 번째 세부 목표에 해당한다.

\subsection{연구 목적 범위의 명확화}
준비 단계 초기에는 ``진단(diagnosis)''과 ``스크리닝(screening)'' 사이의 목적 혼선이 있었으나, 2026년 5월 18일 임상 자문 회의에서 연구 목적은 \textbf{최종 진단이 아니라 앱 사용자의 데이터를 기반으로 PCOS 의심 증상을 스크리닝하여 병원 진료를 권고하는 단계}까지로 범위를 한정하기로 정리되었다. 이 범위 한정은 이후 표본 설계·동의 절차·성능 목표·규제 부담 전반에 영향을 미치는 핵심 결정이었다.

\subsection{보고서의 목적과 성격}
본 보고서는 \textbf{확정된 연구 결과를 보고하는 것이 아니라}, 본격 임상연구 착수 이전에 수행한 \textbf{준비 단계의 현장 활동 경과를 체계적으로 기록}하는 것을 목적으로 한다. 즉, ``어떤 결론이 났는가''가 아니라 ``어떤 데이터를 어떤 방법으로 수집할지 결정하기 위해 어떤 근거를 조사했고, 현장 이해관계자와 무엇을 협의했으며, 무엇이 결정되고 무엇이 미결로 남았는가''를 다룬다. 논문 형식(IMRaD)을 차용하되, 결과 절을 (i)~선행연구 조사, (ii)~프로토콜 설계, (iii)~현장 협의 경과의 세 부분으로 구성하였다.

% =====================================================================
\section{연구 준비 방법 (Methods)}
본 준비 단계는 \textbf{문헌 기반 근거 조사}와 \textbf{현장 기반 이해관계자 협의}를 병행하는 이원 구조로 진행되었다.

\subsection{문헌 조사 방법}
\textbf{탐색 대상 데이터베이스.} PubMed, OpenAlex, Crossref, Semantic Scholar, arXiv, IEEE Xplore, JMIR 계열, Frontiers, Nature/Springer, MDPI, Lancet Digital Health 및 국내 학술 DB(RISS, KCI, DBpia, KoreaScience, KISTI ScienceON).

\smallskip
\noindent\textbf{탐색 절차.} 주제별로 다단계 병렬 탐색을 수행하고, 각 단계 산출물을 누적·합성하였다. 탐색은 다섯 개의 작업공간(workspace)으로 구획하여 진행하였다(표~\ref{tab:workspace}).

\begin{table}[h]
\centering
\footnotesize
\renewcommand{\arraystretch}{1.25}
\caption{문헌 조사 작업공간 구성}
\label{tab:workspace}
\begin{tabularx}{\textwidth}{@{}p{1.9cm} L p{4.0cm}@{}}
\toprule
\textbf{작업공간} & \textbf{조사 주제} & \textbf{핵심 산출물} \\
\midrule
workspace1 & 자궁내막증/PCOS 디지털 바이오마커 일반 + 스마트폰 카메라 바이오마커 & 문헌 리뷰, 바이오마커 카탈로그 \\
workspace2 & PCOS·자궁내막증 특화 카메라 바이오마커 & 문헌 리뷰, 연구설계서 \\
workspace3 & 얼굴 영상·음성 기반 질병 예측 (얼굴 50편 + 음성 43편) & 얼굴·음성 문헌, 합성 보고서 \\
workspace4 & 얼굴·음성 데이터 수집 UX 방법론 + 음성 특징 추출 & UX 합성, 음성 특징 추출 \\
workspace5 & PCOS 자가수집 바이오마커 + 일상 환경(In-the-Wild) 질병 예측 & 합성 프로토콜, 일상환경 프로토콜 \\
\bottomrule
\end{tabularx}
\end{table}

\noindent\textbf{합성 원칙.} 각 문헌의 바이오마커·신체 부위·수집 방법·모델·성능 지표를 정형화하여 추출하고, PCOS Rotterdam 진단 기준과 매핑하여 등급화하였다.

\smallskip
\noindent\textbf{참고문헌 실존성 검증(할루시네이션 검사).} 인용 문헌은 웹 검색(PubMed·PMC·DOI·arXiv 직접 조회)으로 실존 여부를 검증하고, 제목·저자·저널·DOI·표본·성능 수치의 일치 여부를 검증(O)·부분 오류($\triangle$)·미확인(?)·할루시네이션 의심($\times$)으로 분류하였다. 본 보고서 작성 시 헤드라인 문헌 10편을 추가로 재검증하였다(부록~\ref{app:halluc} 참조).

\subsection{현장 협의 방법}
본격 연구의 데이터 수집은 단일 기관으로 완결되지 않고, \textbf{앱 개발사(데이터 수집 채널)·임상 기관(Ground Truth 및 시료)·연구소(분석)}의 협업을 전제로 한다. 따라서 준비 단계의 핵심은 이해관계자 간 합의 형성이었으며, 2026년 4{--}6월에 걸쳐 3회의 현장 협의를 수행하였다(표~\ref{tab:meetings}).

\begin{table}[h]
\centering
\footnotesize
\renewcommand{\arraystretch}{1.25}
\caption{다기관 현장 협의 개요}
\label{tab:meetings}
\begin{tabularx}{\textwidth}{@{}p{0.8cm} p{2.0cm} L p{3.6cm}@{}}
\toprule
\textbf{회차} & \textbf{일시} & \textbf{주요 참석} & \textbf{핵심 안건} \\
\midrule
1차 & 2026-04-27 & 해피문데이(김도진 대표 등), KIST(문경률·류강현·정희은·신홍철·조수민) & 수집 데이터 항목 선정, UX·프라이버시, 질환 타겟팅, 역할 분담 \\
2차 & 2026-05-18 & 해피문데이(김도진), KIST 다수, 고려대병원 박현태 교수 & 연구 목적 재정의, 데이터 통합 전략, 표본 규모, IRB \\
3차 & 2026-06-01 & KIST 다수, 고려대병원 박현태 교수 & 예산 현실화, R\&R 조율, N수 조정, 지식재산권 \\
\bottomrule
\end{tabularx}
\end{table}

% =====================================================================
\section{결과 (I): 선행연구 조사}

\subsection{얼굴 영상 기반 질병 예측 선행연구}
얼굴 영상 분석은 ``한 장의 정지 이미지로 충분한 진단이 가능한 영역''과 ``비디오·시간 동역학이 필요한 영역''으로 양분된다.

\begin{itemize}
  \item \textbf{피부/여드름/호르몬:} 스마트폰 셀피에서 여드름 병변 객체 검출 + 중증도 등급화가 가장 활발하다. AcneDet(Huynh 2022)는 Faster R-CNN + LightGBM으로 IGA 정확도 0.85를 보고하였다~\cite{ref1}. 안드로겐성 여드름은 PCOS의 임상적 주요 징후이나, 여드름$\leftrightarrow$PCOS 직접 연계 모델은 아직 부재하다.
  \item \textbf{눈/결막:} 결막 창백(빈혈), 공막 황달(BiliScreen 계열), 갑상선안병증 등에서 강력한 성능이 보고되었다.
  \item \textbf{표정/얼굴 근육:} 파킨슨 hypomimia, 우울, 뇌졸중/안면마비 검출에 표정 비디오·Action Unit 분석이 사용된다.
  \item \textbf{얼굴 형태/랜드마크:} DeepGestalt(Gurovich 2019, \emph{Nature Medicine})는 200여 개 유전증후군을 Top-10 정확도 91\%로 분류하여 얼굴 형태학 AI의 기념비적 성과로 평가되며, Face2Gene으로 상용화되었다~\cite{ref2}. 말단비대증·쿠싱증후군 등 내분비 질환에서도 전문의를 능가하는 성능이 보고되었다.
  \item \textbf{생물학적 나이/심혈관:} FaceAge(Bontempi 2025, \emph{Lancet Digital Health})는 얼굴 사진으로 생물학적 나이를 추정하여 암 환자 예후 예측을 보강하였다~\cite{ref3}.
\end{itemize}

\subsection{음성 기반 질병 예측 선행연구}
음성 바이오마커는 신경계·정신건강·호흡기·내분비 전반에서 연구되어 왔다.

\begin{itemize}
  \item \textbf{표준 음향 특징:} 기본주파수(F0), jitter, shimmer, HNR, MFCC, 포먼트(F1{--}F4)와 표준화 특징셋(eGeMAPS)이 핵심이다.
  \item \textbf{PCOS·내분비:} Aydin 등은 PCOS 여성에서 안드로겐 과다 $\rightarrow$ 갑상피열근 비대 $\rightarrow$ F0 감소·발성 변화 기전을 일관되게 보고하였다(통계적 음성 변화는 입증, 단 ML 분류 정확도는 미보고)~\cite{ref4}. 갑상선 기능 이상도 음성 변화와 연관된다.
  \item \textbf{표준화 부재:} 음성 수집은 발화 과제(지속 모음 /a/, 음절 반복, 문장 읽기, 자유 발화)의 \textbf{표준화 부재}가 메타분석의 최대 장벽이다.
  \item \textbf{한국어 데이터의 빈곤:} PCOS·갑상선·당뇨 등의 한국어 임상 음성 데이터셋은 사실상 전무하여, 본 과제의 한국 코호트 구축이 기여할 핵심 공백으로 식별되었다.
\end{itemize}

\subsection{자가수집 가능 PCOS 바이오마커 카탈로그}
선행연구를 PCOS Rotterdam 진단 기준(고안드로겐증 / 희발·무배란 / 다낭성난소 형태)과 매핑하여, 스마트폰 자가수집 가능성과 PCOS 연관성에 따라 9종을 등급화하였다(표~\ref{tab:catalog}).

\begin{table}[h]
\centering
\footnotesize
\renewcommand{\arraystretch}{1.2}
\caption{스마트폰 자가수집 가능 PCOS 표현형 바이오마커 카탈로그}
\label{tab:catalog}
\begin{tabularx}{\textwidth}{@{}c L p{2.3cm} p{1.6cm} p{3.0cm} p{2.6cm}@{}}
\toprule
\textbf{등급} & \textbf{바이오마커} & \textbf{부위} & \textbf{자가수집성} & \textbf{성능} & \textbf{근거} \\
\midrule
A & 흑색가시세포증(AN) 색상 & 목 뒷부분 & 높음 & AUC 0.854 (n=227) & ANcam (Dhanoo 2024)~\cite{ref5} \\
A & 여드름 IGA 등급 & 얼굴 3각도 & 높음 & IGA 정확도 0.85 & AcneDet (Huynh 2022)~\cite{ref1} \\
A & 다모증 mFG 9부위 & 9개 신체부위 & 중간 & 일치도 0.89, $\kappa$=0.75 & Oliveira 2023~\cite{ref6} \\
B & rPPG 심박·HRV & 얼굴 영상 30{--}60초 & 높음 & SOTA on-device (n=39) & MobilePhys (Liu 2022)~\cite{ref7} \\
B & 얼굴/전신 BMI & 셀피/전신 사진 & 높음/중간 & MAE 1.04 BMI & 복수 연구 \\
B & 여성형 안드로겐성 탈모 & 정수리·헤어라인 & 중간 & 의사 평가 94\% 일치 & MDhair (Bhardwaj 2025)~\cite{ref8} \\
B & 종단 피부 모니터링 & 표준 4방향 & 중간 & 6개월 추적 & SkinTracker (Jin 2023)~\cite{ref9} \\
C & 공막 색·혈관 패턴 & 안구 8방향 & 낮음 & AUC 0.979 (n=721) & Lv 2022~\cite{ref10} \\
\bottomrule
\end{tabularx}
\end{table}

\paragraph{핵심 통찰.} 자가수집 스마트폰으로 PCOS 3대 진단 기준 중 \textbf{고안드로겐증(다모증+여드름+탈모)을 강하게 평가}할 수 있으며, 인슐린저항성 표현형(AN·BMI)은 보조 진단·중증도 추정에 활용 가능하다. 다낭성난소 형태(초음파)는 자가수집 불가이다. \textbf{PCOS의 4대 표현형을 통합 자가수집하는 연구는 사실상 부재}하여, 이것이 본 연구의 차별점으로 도출되었다.

\subsection{데이터 수집 UX 및 순응도 방법론}
장기 자가수집의 핵심 난제는 \textbf{데이터 품질}과 \textbf{참여자 순응도}이다.
\begin{itemize}
  \item \textbf{캡처 시점 품질 게이트:} Vodrahalli 등(2023, \emph{JAMA Dermatology})은 촬영 시점에 흐림·조명·구도를 실시간 판정하는 AI 지원 도구가 원격 진료 사진 품질을 향상시킴을 입증하였다~\cite{ref11}.
  \item \textbf{지시 없는 자가촬영의 활용성:} Ali 등(2026, \emph{JMIR Dermatology})은 별도 지시 없이 환자가 촬영한 사진의 \textbf{94.6\%가 임상 결정에 충분한 품질}임을 보고하였다~\cite{ref12}.
  \item \textbf{대규모 디지털 코호트의 순응도:} Apple Women's Health Study(Mahalingaiah 2022, \emph{AJOG})는 짧은 자가보고 설문으로 장기 순응도를 유지하며 풍부한 메타데이터를 수집하는 모범을 제시하였다~\cite{ref13}.
  \item \textbf{순응도 설계:} 짧은 온보딩, 자가 모니터링 시각화, 단계적 인센티브가 6개월 retention을 좌우한다(Pratap 2019~\cite{ref14}; EMA 메타분석 Wrzus \& Neubauer 2023~\cite{ref15}).
\end{itemize}

% =====================================================================
\section{결과 (II): 데이터 수집 프로토콜 설계}
선행연구 조사를 바탕으로, 두 가지 상반된 수집 패러다임의 프로토콜을 설계하였다.

\subsection{통제 환경 프로토콜 (v1.0)}
이상적 조건을 사전에 달성하도록 요구하는 엄격 프로토콜로, 6개 촬영 프로토콜(① 얼굴 정·측면, ② 목·겨드랑이 AN, ③ 다모증 9부위, ④ 두피·모발, ⑤ rPPG 영상, ⑥ 전신 BMI)에 대해 거리·각도·조명·시간대·빈도·품질 기준의 구체적 수치를 제시하였다. 통제 조건이 보장되는 서브그룹 비교나 임상 검증 방문 시 활용한다.

\subsection{일상 환경 프로토콜 (v2.0, In-the-Wild) --- 현재 방향}
사용자가 집·사무실·카페 등 \textbf{통제 불가 환경}에서 촬영함을 전제로, \textbf{노이즈를 제거하는 대신 측정·기록·보정}하는 패러다임으로 전환하였다(표~\ref{tab:paradigm}).

\begin{table}[h]
\centering
\footnotesize
\renewcommand{\arraystretch}{1.25}
\caption{수집 패러다임 비교}
\label{tab:paradigm}
\begin{tabularx}{\textwidth}{@{}L L@{}}
\toprule
\textbf{전통 접근 (v1.0)} & \textbf{일상 환경 접근 (v2.0)} \\
\midrule
노이즈 조건 사전 금지 & 노이즈 발생 여부 측정·기록 \\
조건 미충족 시 수집 거부 & 전량 수집 + 품질 등급 A/B/C 자동 부여 \\
동질 조건으로 일반화 & 환경 다양성이 모델 강건성 향상 \\
사전 통제 & 사후 알고리즘 보정 \\
\bottomrule
\end{tabularx}
\end{table}

핵심 설계 요소는 다음과 같다.
\begin{enumerate}
  \item \textbf{절대 거부 조건 최소화(5개만):} 얼굴 미검출, 극단적 과노출, 극단적 블러, 얼굴 크기 부족, (rPPG) 영상 길이 부족. 그 외는 모두 수집한다.
  \item \textbf{품질 메타데이터 자동 수집:} 촬영 시각, 조도, 색온도, 블러 지수, 얼굴 자세, 이동 여부, GPS 환경 유형 등을 분석 단계의 보정 공변량으로 활용한다.
  \item \textbf{참여자 자가보고(10초 이내):} 메이크업 여부, 월경 주기 일차, 직전 활동, 기분·스트레스, 카페인 섭취.
  \item \textbf{참여자 지시 3개만:} 얼굴을 크게, 촬영 중 정지, 표시등 색 확인. 엄격한 금지 목록 대신 ``할 수 있을 때 이렇게 하면 더 좋다'' 수준의 권장만 제공한다.
  \item \textbf{반복 측정에 의한 노이즈 평균화:} 단일 측정의 노이즈보다 누적 측정의 추세가 신뢰성이 높다.
\end{enumerate}

이 설계는 Ali 등(2026)의 ``지시 없는 자가촬영 94.6\% 활용 가능'' 근거~\cite{ref12}와 MobilePhys의 모션 보정~\cite{ref7}, Vodrahalli 등의 실시간 피드백~\cite{ref11}에 기반하며, 생태학적 타당성(앱이 목표하는 일상 환경에서의 예측력)을 우선한다.

\subsection{연구 설계 초안}
준비 단계에서 도출한 연구 설계 초안의 핵심은 다음과 같다(현장 협의에 따라 조정 중).
\begin{itemize}
  \item \textbf{주 가설(H1):} 자가촬영 4대 표현형(여드름·다모증·AN·탈모) 통합 AI 모델이 최고 단일 표현형 모델 대비 PCOS 진단 AUC를 $\geq 0.05$ 향상시킨다.
  \item \textbf{연구 유형:} 다기관 전향적 사례-대조 코호트 + 6개월 종단 자가수집.
  \item \textbf{표본(초안):} PCOS 120 + 대조 120 = 240명(탈락률 50\% 가정 시 모집 목표 480명). \emph{단, 현장 협의에서 표본 규모는 재논의되었다(5절 참조).}
  \item \textbf{개인정보 보호:} 국내 서버 AES-256 암호화, 식별자 분리, 1-탭 철회권, 제3자 공유 금지.
\end{itemize}

% =====================================================================
\section{결과 (III): 다기관 현장 협의 경과}
준비 단계의 핵심 활동은 데이터를 실제로 수집·확보하기 위한 이해관계자 합의 형성이었다. 3회 협의의 경과를 시간순으로 정리한다.

\subsection{1차 협의 (2026-04-27): 데이터 항목 선정과 역할 분담}
\begin{itemize}
  \item \textbf{멀티모달 수집 제안:} 설문(기존 해피문데이 앱 데이터: 월경주기·양·통증·증상 등) + 영상(얼굴·손) + 음성(표준 문장 약 10초, 예: 애국가 따라읽기, 파킨슨 연구 선례 활용)을 통합 수집.
  \item \textbf{UX·프라이버시:} 측정 시간 3{--}5분 이내, 간결한 프로세스로 탈락률 최소화. 얼굴 데이터 전송의 심리적 거부감을 줄이기 위해 \textbf{선글라스 착용·키포인트(keypoint) 추출 저장} 방식을 검토.
  \item \textbf{질환 타겟팅:} PCOS는 선행연구가 존재하여 메인 타겟, 자궁내막증은 의학적 근거 부족으로 박현태 교수 자문을 통해 구체화 예정.
  \item \textbf{표본·기간:} 초기 100{--}200명으로 시작하여 전략 수정, 3년 종단 추적. 일반군 vs PCOS 환자군 분류.
  \item \textbf{역할 분담:} \textbf{앱 개발·데이터 수집 UI는 해피문데이, 데이터 분석·바이오마커 발굴은 KIST.}
  \item \textbf{핵심 쟁점:} 데이터 퀄리티(로우데이터 확보) vs 사용자 편의성(탈락·프라이버시)의 전략적 조율.
\end{itemize}

\subsection{2차 협의 (2026-05-18): 연구 목적 재정의와 데이터 통합 전략}
고려대병원 박현태 교수 참여.
\begin{itemize}
  \item \textbf{목적 재정의:} 연구 목적을 \textbf{진단이 아닌 스크리닝(병원 진료 권고)}으로 명확히 한정(소통 오류 수정).
  \item \textbf{모달리티 확정:} 얼굴(영상·사진)·음성·설문으로 수집. \textbf{웨어러블 센서는 접근성 제한으로 제외}하고 앱 고도화에 집중.
  \item \textbf{데이터 통합 전략:} 디지털 바이오마커(앱) + 멀티오믹스(혈액 시료) + EMR(전자의료기록)의 \textbf{3원 통합 분석}을 본 연구의 강점으로 설정. 멀티오믹스·EMR은 삼성서울병원·고려대병원에서 확보.
  \item \textbf{표본 규모 합의:} \textbf{총 900{--}1{,}500명} 목표(예: PCOS 300, 자궁내막증 400, 정상 300), 우선 100명 규모에서 단계적 확장. 앱 내 쿠폰·유료 서비스 연계로 모집.
  \item \textbf{임상 자문(박현태):} PCOS는 스펙트럼이 넓으므로 \textbf{확진 여부·중증도·약물 복용·치료 진행 여부(Ground Truth)를 명확히 기록}해야 함. 장기 추적 관찰이 유의미.
  \item \textbf{IRB 우려:} 디지털 바이오마커와 멀티오믹스 동시 연계 시, 기존 병원 IRB에 미포함되었을 가능성이 높아 \textbf{재심의/변경 신청 필요(약 2개월 소요 예상).}
  \item \textbf{개발 난이도(김도진):} 영상·음성 수집 기능 개발 난이도가 높으며, 지속 입력을 유도하는 프로토콜 설계가 관건. 예산·기간에 맞춘 우선순위 정리 필요.
\end{itemize}

\subsection{3차 협의 (2026-06-01): 예산 현실화와 역할·권리 조율}
\begin{itemize}
  \item \textbf{예산 현실화:} 약 5억 원 요구는 수용 곤란, \textbf{현재 책정 예산(약 1억 원) 내 진행}으로 정리. 연구비 부족 시 \textbf{얼굴·음성 취득 앱을 KIST가 직접 개발하여 연동}하는 방안 제안.
  \item \textbf{두 질환 동일 파라미터 우려(이지은):} PCOS와 자궁내막증을 \textbf{동일 파라미터로 동시 검출하는 것의 임상적 타당성에 의문}. 현실적으로 \textbf{PCOS에 집중}하고, 자궁내막증은 병원 내부 EMR 활용이 유리할 수 있음.
  \item \textbf{얼굴 데이터 거부감(박현태):} 화장을 지우고 촬영하는 부담 등으로 얼굴 데이터 수집에 회의적 시각.
  \item \textbf{N수 재조정:} 예산·난이도 고려, 각 군 100명 수준 언급되었으나 \textbf{최종 결정 보류.} 확진 그룹은 기존 연구 기준 최소 200{--}250명 필요(박현태).
  \item \textbf{IRB·소유권:} 오믹스와 디지털 바이오마커를 분리 진행하기 어려워 \textbf{선행(preliminary) 연구 형태로 심사} 전망. 지식재산권은 혁신법상 \textbf{공동기관 소유} 가능(옥명렬).
  \item \textbf{자궁내막증 서브 과제(박현태):} 산부인과 초음파·MRI 데이터(차트·초음파·MRI 수백 건)를 자궁내막증 연구에 서브로 연계하는 방안 제안.
\end{itemize}

\subsection{현장 협의 종합: 합의 사항과 미결 쟁점}
\begin{table}[h]
\centering
\footnotesize
\renewcommand{\arraystretch}{1.3}
\caption{현장 협의 종합 --- 합의(수렴) 대 미결(쟁점)}
\label{tab:summary}
\begin{tabularx}{\textwidth}{@{}>{\bfseries}p{2.0cm} L@{}}
\toprule
합의(수렴) & ① 목적 = 스크리닝 / ② 모달리티 = 얼굴·음성·설문(웨어러블 제외) / ③ 3원 통합(디지털 바이오마커+오믹스+EMR) / ④ 역할 = 해피문데이 수집·KIST 분석 / ⑤ PCOS 우선·단계적 확장 / ⑥ 장기 추적 지향 \\
\midrule
미결(쟁점) & ① 최종 표본 규모(100 vs 1{,}000) / ② 예산 배분(1억 vs 5억) / ③ 얼굴 데이터 수용도·익명화 방식 / ④ 자궁내막증의 별도 경로(EMR/초음파) 처리 / ⑤ IRB 재심의 범위·일정 / ⑥ 지식재산권·성과 배분 / ⑦ 음성 발화 프로토콜 표준화 \\
\bottomrule
\end{tabularx}
\end{table}

% =====================================================================
\section{고찰 (Discussion)}

\subsection{본 준비 단계의 의의}
본 단계는 \textbf{확정적 결론을 산출하기 위한 것이 아니라, 근거 기반 의사결정을 위한 토대를 마련하는 과정}이었다. 선행연구 조사는 ``무엇을 측정할 수 있는가''(자가수집 가능 바이오마커 9종)를, 프로토콜 설계는 ``어떻게 측정할 것인가''(일상 환경 측정·보정 패러다임)를, 현장 협의는 ``현실적으로 무엇을 수집할 수 있는가''(이해관계자 제약)를 각각 규명하였다. 세 축이 교차하는 지점에서 실현 가능한 연구 방향이 수렴되고 있다.

\subsection{근거와 현장의 긴장 관계}
흥미롭게도 \textbf{문헌 근거가 지지하는 방향과 현장 제약이 충돌하는 지점}이 다수 확인되었다.
\begin{itemize}
  \item \textbf{로우데이터 vs 수용도:} 분석 정확도를 위해서는 로우데이터·얼굴 영상이 유리하나(문헌 지지), 현장에서는 얼굴 데이터의 심리적 거부감과 개발 난이도가 큰 장벽이다. $\rightarrow$ 키포인트 추출·익명화·일상 환경 프로토콜(거부 최소화)이 절충안으로 도출되었다.
  \item \textbf{표본 규모:} 통합 모델 검증에는 수백 명이 필요하나(초안 480명), 예산 한계로 100명 단계 시작이 현실적이다. $\rightarrow$ 단계적 확장 설계로 조정.
  \item \textbf{두 질환 동시 추진:} 과제 목표는 PCOS·자궁내막증 모두이나, 자궁내막증은 디지털 바이오마커 근거가 부족하다(문헌·임상 자문 일치). $\rightarrow$ PCOS 우선, 자궁내막증은 EMR/초음파 별도 경로로 분기.
\end{itemize}

\subsection{연구 공백과 차별점}
선행연구 조사 결과, \textbf{PCOS + 카메라 자가수집 + 일상 환경 + 다중 표현형 통합}의 교차 영역은 사실상 백지 상태이다. Apple Women's Health Study(10만 명+)조차 카메라를 활용하지 않았다~\cite{ref13}. 또한 PCOS·갑상선 등의 한국어 음성·얼굴 임상 데이터셋이 전무하여, 한국 코호트 구축 자체가 학술적 기여가 된다.

\subsection{한계}
(1)~본 보고서는 데이터 수집 착수 이전 단계의 기록으로, 실제 수집·분석 결과는 포함하지 않는다. (2)~표본·예산·IRB 등 핵심 변수가 미확정이어서 연구 설계는 잠정적이다. (3)~자가촬영 mFG의 임상 일치도, rPPG-PCOS 직접 연관성, 한국 인구 적용 가능성 등은 본 연구가 처음 검증해야 할 가정이다. (4)~자궁내막증의 디지털 바이오마커 경로는 본 단계에서 충분히 규명되지 못했다.

% =====================================================================
\section{결론 및 향후 과제}
본 현장연구는 스마트폰 기반 얼굴·음성 디지털 바이오마커로 PCOS를 스크리닝하는 임상연구의 \textbf{착수 준비 활동}을 수행하였다. 선행연구 조사로 자가수집 가능 바이오마커 9종을 등급화하고, 일상 환경 수집·보정 프로토콜을 설계하였으며, 3회의 다기관 현장 협의로 수집 모달리티·역할 분담·PCOS 우선 추진에 대한 합의를 이루었다. 다만 \textbf{표본 규모·예산·얼굴 데이터 수용도·IRB·지식재산권은 미결 쟁점으로 남아}, 본 단계는 명확한 단일 결론보다는 \textbf{의사결정의 근거와 선택지를 정련하는 과정}으로서의 성격을 가진다.

\paragraph{향후 과제.}
\begin{enumerate}
  \item \textbf{해피문데이 재협의:} 예산·역할·앱 개발 범위·데이터 소유권의 오해 해소(문경률·류강현 주도).
  \item \textbf{IRB 전략 수립:} 오믹스+디지털 바이오마커 통합을 선행 연구 형태로 심의, 재심의 일정 확보.
  \item \textbf{음성 발화 프로토콜 표준화:} 선행 논문 기반 표준 발화 과제(문장·시간·빈도) 확정.
  \item \textbf{자궁내막증 경로 분리 설계:} EMR·초음파/MRI 기반 별도(서브) 과제로 구체화.
  \item \textbf{표본 크기·통계 분석 계획 확정:} 단계적 확장(100명 $\rightarrow$ 1{,}000명) 설계의 통계적 검정력 재계산.
  \item \textbf{삼성서울병원 협의:} 멀티오믹스·시료 제공 및 IRB 정합성 확보.
\end{enumerate}

% =====================================================================
\section*{부록 A. 참고문헌 할루시네이션 검증 결과}
\addcontentsline{toc}{section}{부록 A. 참고문헌 할루시네이션 검증 결과}
\label{app:halluc}

본 보고서에 인용한 핵심 헤드라인 문헌에 대해, 작성 시점(2026-06-16)에 웹 검색(PubMed·PMC·DOI·arXiv·학회 DB 직접 조회)으로 실존 여부 및 서지정보 일치를 재검증하였다.

\begin{table}[h]
\centering
\footnotesize
\renewcommand{\arraystretch}{1.2}
\caption{헤드라인 문헌 재검증 결과 (O=일치, $\triangle$=부분오류)}
\begin{tabularx}{\textwidth}{@{}c L p{4.2cm} c@{}}
\toprule
\textbf{\#} & \textbf{문헌} & \textbf{검증 항목} & \textbf{결과} \\
\midrule
\cite{ref1} & Huynh 2022 AcneDet, \emph{Diagnostics} & 제목·저자·DOI·표본(1{,}572장) & O \\
\cite{ref2} & Gurovich 2019 DeepGestalt, \emph{Nat Med} & 제목·저자·권/페이지·DOI·PMID & O \\
\cite{ref4} & Aydin PCOS voice (복수 논문) & 실존·기전(안드로겐$\rightarrow$F0 감소) & O \\
\cite{ref5} & Dhanoo 2024 ANcam, \emph{Diabetes Spectrum} & 제목·DOI·표본·성능(AUC 0.854) & $\triangle$ \\
\cite{ref6} & Oliveira 2023 mFG, \emph{Arch Dermatol Res} & 제목·저자·PMID·일치도(0.89) & O \\
\cite{ref7} & Liu 2022 MobilePhys, \emph{ACM IMWUT} & 제목·저자·DOI·arXiv & O \\
\cite{ref10} & Lv 2022 PCOS 공막, \emph{Front Endocrinol} & 제목·저자·DOI·PMID·성능(AUC 0.979) & O \\
\cite{ref11} & Vodrahalli 2023, \emph{JAMA Dermatology} & 제목·PMC·표본(98명/357장) & O \\
\cite{ref12} & Ali 2026, \emph{JMIR Dermatology} & 제목·저자·DOI·PMC·수치(94.6\%) & O \\
\cite{ref13} & Mahalingaiah 2022 Apple WHS, \emph{AJOG} & 제목·저자·권/페이지·PMC & O \\
\bottomrule
\end{tabularx}
\end{table}

\textbf{결론.} 재검증한 헤드라인 문헌 10건 모두 \textbf{실존이 확인}되었으며, 9건은 서지정보가 정확하게 일치하였다. 1건(ANcam)은 저자명 오류가 발견되어 정정하였다(할루시네이션 의심이 아닌 서지 표기 오류).

\paragraph{발견된 오류 및 정정 --- ANcam (Dhanoo 2024).}
기존 합성보고서 표기 ``Dhanoo D, Greene Z, Berbesi-Fernandez D, et al.''는 검증된 실제 저자 \textbf{Andrew S. Dhanoo, Sterling K. Ramroach, Felicia Hill-Briggs, Brian N. Cockburn}와 불일치하여 정정하였다. 논문 자체·제목·DOI(10.2337/ds23-0042)·성능(AUC 0.854, 민감도 81.1\%, 특이도 70.3\%, n=227)은 모두 정확하다.

\paragraph{본문 인용에서 제외한 항목.}
기존 검증에서 할루시네이션 의심($\times$)으로 분류된 arXiv:1907.07901(셀피 여드름, 저자 정보 전체 오류)과 미검증(?)인 Cao et al.\ 2025 ECEESPE2025 P804(미발표 학술대회 abstract)은 보수적 원칙에 따라 \textbf{본문 인용에서 제외}하였다.

\paragraph{검증 방법 및 한계.}
WebSearch를 통한 PubMed/PMC/DOI/arXiv/학회 DB 직접 조회로 제목·저자·저널·권호·DOI·PMID·표본·성능 수치를 교차 대조하였다. 한계로는 (1)~US 기반 검색으로 일부 국내 학술 DB·비색인 자료의 직접 확인 제한, (2)~미발표 abstract·그레이 문헌의 보수적 제외, (3)~본 부록은 헤드라인 10편 재검증이며 전체 문헌의 상세 검증은 각 작업공간 검증보고서를 따른다는 점이 있다.

% =====================================================================
\begin{thebibliography}{99}

\bibitem{ref1} Huynh QT, Nguyen PH, Le HX, et al. Automatic Acne Object Detection and Acne Severity Grading Using Smartphone Images and Artificial Intelligence. \emph{Diagnostics}. 2022;12(8):1879. PMID:36010229. \href{https://doi.org/10.3390/diagnostics12081879}{doi:10.3390/diagnostics12081879}.

\bibitem{ref2} Gurovich Y, Hanani Y, Bar O, et al. Identifying facial phenotypes of genetic disorders using deep learning. \emph{Nature Medicine}. 2019;25(1):60--64. PMID:30617323. \href{https://doi.org/10.1038/s41591-018-0279-0}{doi:10.1038/s41591-018-0279-0}.

\bibitem{ref3} Bontempi D, et al. FaceAge: deep learning estimation of biological age from face photographs for survival prediction. \emph{Lancet Digital Health}. 2025. (Code: GitHub AIM-Harvard/FaceAge).

\bibitem{ref4} Aydin K, et al. Voice characteristics associated with polycystic ovary syndrome. \emph{The Laryngoscope}. 2016. PMID:26700739. (관련: Vocal Changes in Patients With PCOS. \emph{Journal of Voice}. 2010, PMID:20537860; Aydin 2024, \emph{Egyptian Journal of Otolaryngology}, \href{https://doi.org/10.1186/s43163-024-00659-5}{doi:10.1186/s43163-024-00659-5}.)

\bibitem{ref5} Dhanoo AS, Ramroach SK, Hill-Briggs F, Cockburn BN. Grading Acanthosis Nigricans Using a Smartphone and Color Analysis (ANcam). \emph{Diabetes Spectrum}. 2024;37(2):139--148. PMID:38756432. \href{https://doi.org/10.2337/ds23-0042}{doi:10.2337/ds23-0042}.

\bibitem{ref6} Oliveira TF, Rocha ALL, Reis FM, C\^{a}ndido AL, Premaor MO, Comim FV. Comparison of image-based modified Ferriman-Gallwey score evaluation with in-person evaluation: an alternative method for hirsutism diagnosis. \emph{Archives of Dermatological Research}. 2023;315(7):1949--1955. PMID:36508021. \href{https://doi.org/10.1007/s00403-022-02495-0}{doi:10.1007/s00403-022-02495-0}.

\bibitem{ref7} Liu X, Wang Y, Xie S, Zhang X, Ma Z, McDuff D, Patel S. MobilePhys: Personalized Mobile Camera-Based Contactless Physiological Sensing. \emph{Proc.\ ACM IMWUT}. 2022;6(1). \href{https://doi.org/10.1145/3517225}{doi:10.1145/3517225}. arXiv:2201.04039.

\bibitem{ref8} Bhardwaj V, Rodgers N, Harth O, Harth Y. AI-Based Personalization of Treatment Regimen for Hair Loss: A 6-Month Clinical Trial (MDhair). \emph{Journal of Drugs in Dermatology}. 2025. (Article S1545961625P8611X).

\bibitem{ref9} Jin JQ, et al. Development of SkinTracker, an integrated dermatology mobile app and web portal enabling remote clinical research studies. \emph{Frontiers in Digital Health}. 2023;5:1228503. PMID:37744686. \href{https://doi.org/10.3389/fdgth.2023.1228503}{doi:10.3389/fdgth.2023.1228503}.

\bibitem{ref10} Lv W, Song Y, Fu R, et al. Deep Learning Algorithm for Automated Detection of Polycystic Ovary Syndrome Using Scleral Images. \emph{Frontiers in Endocrinology}. 2022;12:789878. PMID:35154003. \href{https://doi.org/10.3389/fendo.2021.789878}{doi:10.3389/fendo.2021.789878}.

\bibitem{ref11} Vodrahalli K, Ko J, Zou J, Daneshjou R, et al. Development and Clinical Evaluation of an Artificial Intelligence Support Tool for Improving Telemedicine Photo Quality. \emph{JAMA Dermatology}. 2023. \href{https://pmc.ncbi.nlm.nih.gov/articles/PMC10018405/}{PMC10018405}.

\bibitem{ref12} Ali Z, Thomsen K, Vestergaard C, Thomsen SF. Assessment of Quality and Utility of Patient-Taken Smartphone Photographs of Atopic Dermatitis: Clinical Survey Study. \emph{JMIR Dermatology}. 2026;9:e72916. \href{https://doi.org/10.2196/72916}{doi:10.2196/72916}. PMC12844854.

\bibitem{ref13} Mahalingaiah S, Fruh V, Rodriguez E, et al. Design and methods of the Apple Women's Health Study: a digital longitudinal cohort study. \emph{American Journal of Obstetrics \& Gynecology}. 2022;226(4):545.e1--545.e29. \href{https://pmc.ncbi.nlm.nih.gov/articles/PMC10518829/}{PMC10518829}.

\bibitem{ref14} Pratap A, et al. Indicators of retention in remote digital health studies. \emph{npj Digital Medicine}. 2019. \href{https://pmc.ncbi.nlm.nih.gov/articles/PMC6483978/}{PMC6483978}.

\bibitem{ref15} Wrzus C, Neubauer AB. Ecological Momentary Assessment: A Meta-Analysis on Designs, Samples, and Compliance Across Research Fields. \emph{Assessment}. 2023. \href{https://doi.org/10.1177/10731911211067538}{doi:10.1177/10731911211067538}. PMC9999286.

\end{thebibliography}

\end{document}
